% Standalone problem document, generated from the adjacent README.md.
% Compile with XeLaTeX. Mathematical content is maintained in Markdown.
\documentclass[11pt,a4paper]{article}
\usepackage[fontsize=10.5pt]{scrextend}
\usepackage[margin=22mm,headheight=15pt,headsep=9mm,footskip=11mm]{geometry}
\usepackage{fontspec}
\setmainfont{texgyrepagella}[Extension=.otf,UprightFont=*-regular,BoldFont=*-bold,ItalicFont=*-italic,BoldItalicFont=*-bolditalic]
\setsansfont{texgyreheros}[Extension=.otf,UprightFont=*-regular,BoldFont=*-bold,ItalicFont=*-italic,BoldItalicFont=*-bolditalic]
\setmonofont{texgyrecursor}[Extension=.otf,UprightFont=*-regular,BoldFont=*-bold,ItalicFont=*-italic,BoldItalicFont=*-bolditalic,Scale=MatchLowercase]
\usepackage{amsmath,amssymb}
\usepackage{unicode-math}
\setmathfont{texgyrepagella-math.otf}
\usepackage{xcolor,microtype,enumitem,needspace,fancyhdr}
\usepackage{bookmark}
\definecolor{ink}{HTML}{17344B}
\definecolor{accent}{HTML}{197D80}
\definecolor{muted}{HTML}{596773}
\hypersetup{colorlinks=true,linkcolor=accent,urlcolor=accent,pdftitle={IE-14 - Sharp
growth for cyclic tridiagonal
matrices},pdfauthor={Open Problems in Numerical Linear Algebra}}
\urlstyle{same}
\setlength{\parindent}{0pt}
\setlength{\parskip}{5pt plus 1pt minus 1pt}
\linespread{1.04}
\setlength{\emergencystretch}{3em}
\widowpenalty=10000
\clubpenalty=10000
\displaywidowpenalty=10000
\allowdisplaybreaks[1]
\setlist{nosep,leftmargin=1.5em}
\providecommand{\tightlist}{\setlength{\itemsep}{0pt}\setlength{\parskip}{0pt}}
\providecommand{\pandocbounded}[1]{#1}
\pagestyle{fancy}
\fancyhf{}
\fancyhead[L]{\sffamily\footnotesize\color{muted}OPEN PROBLEMS IN NUMERICAL LINEAR ALGEBRA}
\fancyhead[R]{\sffamily\bfseries\footnotesize\color{accent}IE-14}
\fancyfoot[L]{\parbox[b]{0.9\textwidth}{\sffamily\fontsize{8}{10}\selectfont\color{muted}Editorial ratings; a bounded search cannot certify that a problem remains unsolved.\\Literature check: 2026-09-22}}
\fancyfoot[R]{\sffamily\footnotesize\color{muted}\thepage}
\renewcommand{\headrulewidth}{0.3pt}
\renewcommand{\headrule}{\hbox to\headwidth{\color{accent}\leaders\hrule height \headrulewidth\hfill}}
\makeatletter
\renewcommand{\section}{\Needspace{5\baselineskip}\@startsection{section}{1}{0pt}{12pt plus 2pt minus 2pt}{4pt}{\sffamily\bfseries\large\color{ink}}}
\renewcommand{\subsection}{\Needspace{4\baselineskip}\@startsection{subsection}{2}{0pt}{9pt plus 1pt minus 1pt}{3pt}{\sffamily\bfseries\normalsize\color{ink}}}
\makeatother
\setcounter{secnumdepth}{0}
\begin{document}
{\sffamily\small\color{muted}Linear systems and elimination\par}
\vspace{3pt}
{\raggedright\hyphenpenalty=10000\exhyphenpenalty=10000\sffamily\bfseries\fontsize{21}{24}\selectfont\color{ink}Sharp
growth for cyclic tridiagonal matrices\par}
\vspace{7pt}
{\sffamily\small\textbf{Difficulty:} challenging\quad\textbf{Importance:} interesting
to specialist\par}
{\sffamily\small\bfseries\color{accent}Status: Lean verified\par}
\vspace{4pt}
{\color{accent}\hrule height 0.8pt}
\vspace{6pt}
\subsection{Lean verification --- 22 September
2026}\label{lean-verification-22-september-2026}

The complete original target is formally proved: for every \(n\ge4\),
the sharp growth factor is \(c_n=F_{n+1}+1\). The upper bound covers
every nonsingular complex cyclic tridiagonal input with both cyclic
corners nonzero, every permitted GEPP tie path, and every active Schur
entry in the original ordering. The explicit rational family attains the
bound through actual row swaps. The nonempty bounded growth set, its
greatest member and its genuine real supremum are all proved.

All four exports passed LeanCert kernel audits, two independent
nonauthor final reviews, and
\href{https://github.com/sidneyholden1/OpenProblemsInNLA/actions/runs/35700510972}{isolated
Linux Comparator/default-kernel verification} with actual rejection and
sandbox controls.
\href{https://github.com/sidneyholden1/OpenProblemsInNLA/blob/main/linear-systems-and-elimination/IE-14/lean/README.md}{Proof
and reviews} ·
\href{https://github.com/sidneyholden1/OpenProblemsInNLA/tree/9a0117aa8ed40a73b285772b3b1c13864782e95d/linear-systems-and-elimination/IE-14/lean}{Immutable
proof revision} ·
\href{https://github.com/sidneyholden1/OpenProblemsInNLA/blob/main/docs/lean/verification-2026-09-22/IE-14/README.md}{Retained
evidence}.

Mathematics: Matthew J. Colbrook; Higham retains original-problem
credit. Formalization: Sidney Holden with OpenAI Codex assistance. No
symmetry, diagonal dominance, real-only restriction or deterministic tie
rule is substituted. No novelty, source-author endorsement or external
human peer review is claimed. The complete original target and earlier
source record remain below.

\subsection{Independently reviewed resolution -
2026-09-11}\label{independently-reviewed-resolution---2026-09-11}

\textbf{Sharp growth classification.} Theorem 1 and Sections 2-4 prove
\(c_n=F_{n+1}+1\) for every \(n\ge4\), with \(F_0=0,F_1=1\). The bound
covers complex cyclic tridiagonal matrices, every active entry and every
permitted GEPP tie path in the original ordering. A rational matrix with
both cyclic corners nonzero attains it at every order.

\textbf{Author:} Matthew J. Colbrook, Department of Applied Mathematics
and Theoretical Physics, University of Cambridge.
\href{https://github.com/ajt60gaibb/OpenProblemsInNLA/blob/main/references/colbrook-recovered-2026-09-11/manuscripts/IE-14.pdf}{Complete
manuscript},
\href{https://github.com/ajt60gaibb/OpenProblemsInNLA/blob/main/references/colbrook-recovered-2026-09-11/verification/reviews/IE-14-review.md}{independent
proof review}, and
\href{https://github.com/ajt60gaibb/OpenProblemsInNLA/blob/main/references/colbrook-recovered-2026-09-11/README.md}{submission
record}. The supplied notes were reconstructed with substantial AI
assistance. This is independent agent verification, not external human
peer review or formal proof-assistant certification; no novelty or
priority claim is made.

The difficulty, importance and rating rationale below are historical
assessments of the original open target. Original statements, references
and dated audits are preserved.

\textbf{Rating rationale:} Challenging reflects an all-orders extremal
problem in which the two corner entries change elimination fill;
specialist impact is a precise stability bound for cyclic tridiagonal
systems.

\newpage
\subsection{Problem statement}\label{problem-statement}

For each \(n\ge4\), let \(\mathcal C_n\) consist of the nonsingular
matrices \(A\in\mathbb C^{n\times n}\) whose only permitted nonzero
positions satisfy

\[|i-j|\le1\quad\text{or}\quad(i,j)\in\{(1,n),(n,1)\},\]

with \(a_{1n}\ne0\) and \(a_{n1}\ne0\). This is the cyclic tridiagonal
pattern called \emph{quasi-tridiagonal} in the source. No symmetry or
diagonal dominance is assumed.

Apply Gaussian elimination with partial pivoting (GEPP) in exact
arithmetic. At each stage select a largest-modulus entry in the active
first column, interchange its row with the active first row, and form
the trailing Schur complement. For an admissible path \(\pi\), let
\(S_1=A,\ldots,S_n\) denote the active matrices. Define

\[\rho(A,\pi)=\frac{\max_{1\le k\le n}\|S_k\|_{\max}}{\|A\|_{\max}},
\qquad \|M\|_{\max}=\max_{i,j}|m_{ij}|,\]

\[c_n=\sup_{A\in\mathcal C_n}\ \sup_{\pi\text{ permitted by GEPP}}\rho(A,\pi).\]

Determine \(c_n\) for every \(n\ge4\), with matching upper bounds and
examples approaching each supremum. All admissible tie choices are
included. The given ordering is fixed before the stated GEPP run;
preliminary row or column reorderings are not permitted. All sizes
constitute one extremal problem.

Ordinary tridiagonal inputs have growth at most two. This question
isolates the effect of the two corner entries on partial-pivoting
stability. Unlike the fixed-bandwidth problem IE-13, the corner offsets
grow with \(n\).

\subsection{Reference}\label{reference}

N. J. Higham,
\href{https://doi.org/10.1137/1.9780898718027}{\emph{Accuracy and
Stability of Numerical Algorithms}, second edition}, SIAM (2002),
Problem 9.15(b), p.~193; Theorem 9.11, p.~173, supplies the ordinary
tridiagonal comparison.

\subsection{Earlier status check ---
2026-09-08}\label{earlier-status-check-2026-09-08}

The book's sparsity pattern was checked directly. Its wording omits the
field; this entry adopts \(\mathbb C\) from the adjacent Theorem 9.11.
Searches for \textit{quasi-tridiagonal growth factor},
\textit{quasi-tridiagonal partial pivoting}, and
\textit{Gaussian elimination cyclic tridiagonal growth} found no
solution of this exact extremal question. A
\href{https://doi.org/10.1007/s10910-017-0761-9}{2017 cyclic-reduction
paper} uses a different matrix pattern and algorithm. Shah--Urschel's
\href{https://arxiv.org/html/2608.19189v4}{August 31, 2026 revision},
Theorems 2.2--2.3, does not determine this extremum. No recent explicit
reaffirmation of the question's openness was found; absence of a located
resolution is the limit of this check.

\subsection{Audit update --- 2026-09-10}\label{audit-update-2026-09-10}

Rechecked Higham's
\href{https://pages.stat.wisc.edu/~bwu62/771/hingham2002.pdf}{Problem
9.15(b)}, p.~193, against the two-corner pattern displayed here.
Searches for cyclic/quasi-tridiagonal pivot growth and later elimination
bounds found no exact solution. The
\href{https://arxiv.org/html/2608.19189v4}{August 2026 general
complete-pivoting results} address a different pivot rule and do not
provide this pattern's sharp partial-pivoting bound; current openness
rests on the historical question and this limited search.
\end{document}
